% Options for packages loaded elsewhere
\PassOptionsToPackage{unicode}{hyperref}
\PassOptionsToPackage{hyphens}{url}
\PassOptionsToPackage{dvipsnames,svgnames,x11names}{xcolor}
\documentclass[
  10pt,
]{ctexart}
\usepackage{xcolor}
\usepackage[margin=2cm]{geometry}
\usepackage{amsmath,amssymb}
\setcounter{secnumdepth}{-\maxdimen} % remove section numbering
\usepackage{iftex}
\ifPDFTeX
  \usepackage[T1]{fontenc}
  \usepackage[utf8]{inputenc}
  \usepackage{textcomp} % provide euro and other symbols
\else % if luatex or xetex
  \usepackage{unicode-math} % this also loads fontspec
  \defaultfontfeatures{Scale=MatchLowercase}
  \defaultfontfeatures[\rmfamily]{Ligatures=TeX,Scale=1}
\fi
\usepackage{lmodern}
\ifPDFTeX\else
  % xetex/luatex font selection
  \setmainfont[]{DejaVu Sans}
\fi
% Use upquote if available, for straight quotes in verbatim environments
\IfFileExists{upquote.sty}{\usepackage{upquote}}{}
\IfFileExists{microtype.sty}{% use microtype if available
  \usepackage[]{microtype}
  \UseMicrotypeSet[protrusion]{basicmath} % disable protrusion for tt fonts
}{}
\makeatletter
\@ifundefined{KOMAClassName}{% if non-KOMA class
  \IfFileExists{parskip.sty}{%
    \usepackage{parskip}
  }{% else
    \setlength{\parindent}{0pt}
    \setlength{\parskip}{6pt plus 2pt minus 1pt}}
}{% if KOMA class
  \KOMAoptions{parskip=half}}
\makeatother
\usepackage{longtable,booktabs,array}
\usepackage{caption}
\captionsetup[table]{skip=6pt}
\newcounter{none} % for unnumbered tables
\usepackage{calc} % for calculating minipage widths
% Correct order of tables after \paragraph or \subparagraph
\usepackage{etoolbox}
\makeatletter
\patchcmd\longtable{\par}{\if@noskipsec\mbox{}\fi\par}{}{}
\makeatother
% Allow footnotes in longtable head/foot
\IfFileExists{footnotehyper.sty}{\usepackage{footnotehyper}}{\usepackage{footnote}}
\makesavenoteenv{longtable}
\setlength{\emergencystretch}{3em} % prevent overfull lines
\providecommand{\tightlist}{%
  \setlength{\itemsep}{0pt}\setlength{\parskip}{0pt}}
\usepackage{bookmark}
\IfFileExists{xurl.sty}{\usepackage{xurl}}{} % add URL line breaks if available
\urlstyle{same}
% fallback for those not using the hyperref driver hyperxmp:
\makeatletter
\@ifundefined{xmpquote}{\newcommand{\xmpquote}[1]{#1}}{}
\makeatother
\hypersetup{
  colorlinks=true,
  linkcolor={Maroon},
  filecolor={Maroon},
  citecolor={Blue},
  urlcolor={Blue},
  pdfcreator={LaTeX via pandoc}}

\author{}
\date{}

\usepackage{pdflscape,fvextra}
\setlength{\tabcolsep}{3pt}
\begin{document}

\section{LLM
知识污染探针（TRANSCRIPT，v1）}\label{llm-ux77e5ux8bc6ux6c61ux67d3ux63a2ux9488transcriptv1}

运行日期：2026-09-28。模型：\texttt{deepseek-flash}（DeepSeek beta
端点，strict function calling，\texttt{temperature=0}）。
运行器：\texttt{evals/\allowbreak{}contamination\_\allowbreak{}probe.\allowbreak{}py}；聚合结果：\texttt{benchmark/\allowbreak{}results/\allowbreak{}contamination\_\allowbreak{}probe\_\allowbreak{}v1.\allowbreak{}json}（只含聚合数字、样本哈希、调用计数和成本估计，不含原始
prompt/response）。 每次模型请求都由 \texttt{TraceRecorder} 单独写入带
SHA-256 链的 trace（位于被忽略的
\texttt{artifacts/\allowbreak{}contamination\_\allowbreak{}probe/\allowbreak{}provider\_\allowbreak{}traces/\allowbreak{}}），入口级
trace 位于 \texttt{artifacts/\allowbreak{}contamination\_\allowbreak{}probe/\allowbreak{}traces/\allowbreak{}}。

\subsection{目的}\label{ux76eeux7684}

检验"让 LLM
看到药物/疾病\textbf{名称}"时，它能否凭记忆中的背景知识"找回" TRANSCRIPT
中已知的适应症，也就是标签泄漏。如果能，那么任何让模型看到身份信息的基准测得的只是记忆，不是发现能力。本项目的
benchmark 采用 strict 模式（LLM
永远看不到药物/疾病身份），这个实验用来检验这条规则是否必要。

\subsection{协议}\label{ux534fux8bae}

\begin{enumerate}
\def\labelenumi{\arabic{enumi}.}
\tightlist
\item
  \textbf{ID→名称映射（先查离线来源）}：仓库内的 Repurposing Hub
  文件只有 Broad ID、名称、InChIKey 和 PubChem CID，没有 DrugBank ID 或
  MedGen CUI，仓库里也没有现成对照表，所以改用公开 API：

  \begin{itemize}
  \tightlist
  \item
    DrugBank ID → mychem.info \texttt{drugbank.\allowbreak{}name}（610 个 DB ID）；
  \item
    PubChem CID → PubChem PUG REST \texttt{Title}（3 个 CID）；
  \item
    MedGen CUI → NCBI E-utilities
    \texttt{esearch}/\texttt{esummary}（db=medgen）的 \texttt{title}。
    映射缓存在
    \texttt{artifacts/\allowbreak{}contamination\_\allowbreak{}probe/\allowbreak{}}（被忽略），覆盖率汇总见
    \texttt{mapping\_\allowbreak{}coverage.\allowbreak{}json}。
  \end{itemize}
\item
  \textbf{抽样并冻结}（seed
  20260928）：在药物名和疾病名都已映射的配对中，抽取阳性（1）并以 300
  为上限随机抽样，再随机抽取同样数量的未知（0）配对，另外纳入全部 −1
  配对。样本在\textbf{任何模型调用之前}写成规范化 JSONL 并记录
  SHA-256；之后每一步都会重新校验哈希，哈希不同的样本拒绝覆盖。
\item
  \textbf{三种条件}：

  \begin{itemize}
  \tightlist
  \item
    \textbf{Open-book}：只问 ``Is an approved or established treatment
    for ?''，要求返回 0--1 概率（strict JSON 工具调用）。
  \item
    \textbf{Closed-book}：不提供任何名称或 ID，只提供该配对的 Spearman
    反转分数（items.csv 与 users.csv 在 12,096 个基因上的负 Spearman
    相关），以及它在该药物 151
    个疾病中的排名百分位，问同一个概率问题。运行器在发送前会断言请求中不含名称、ID
    或标签。
  \item
    \textbf{原始反转分数}：直接用 Spearman 反转分数作为排序分数，不经过
    LLM。
  \end{itemize}
\item
  \textbf{指标}：主指标是阳性 vs 未知的 AUC（Mann--Whitney，并列按 0.5
  计），用分层 bootstrap（2000 次，seed 20260928）求 95\% CI。open 与
  closed 的差值用配对分层 bootstrap。阳性 vs −1 只作描述，因为 n=11。
\end{enumerate}

\subsection{结果}\label{ux7ed3ux679c}

映射覆盖率：药物 611/613（未映射 DB00510、DB01361），疾病
150/151（未映射 C0023473）；阳性 399/401、−1 11/11 的两端名称都已映射。
样本：300 个阳性、300 个未知、11 个 −1，共 611 对；SHA-256 为
\texttt{898a59fe00cc1d\allowbreak{}bcdb22a0690773\allowbreak{}7eb723f1b1bbc9\allowbreak{}c36cc188f59ad4\allowbreak{}22eb4e4b}。

{\def\LTcaptype{none} % do not increment counter
\begin{longtable}[]{@{}>{\raggedright\arraybackslash}p{0.1800\textwidth}>{\raggedright\arraybackslash}p{0.1800\textwidth}>{\raggedright\arraybackslash}p{0.1800\textwidth}>{\raggedright\arraybackslash}p{0.1800\textwidth}>{\raggedright\arraybackslash}p{0.1800\textwidth}@{}}
\toprule\noalign{}
条件 & AUC（阳性 vs 未知） & 95\% CI & 阳性 / 未知的中位概率 & 阳性 vs
−1（描述，n=11） \\
\midrule\noalign{}
\endhead
\bottomrule\noalign{}
\endlastfoot
Open-book（看到名称） & \textbf{0.567} & {[}0.526, 0.609{]} & 0.02 /
0.02 & 0.516 {[}0.362, 0.672{]} \\
Closed-book（只看反转数字） & 0.527 & {[}0.482, 0.570{]} & 0.006 / 0.006
& 0.481 {[}0.312, 0.636{]} \\
原始 Spearman 反转分数 & 0.519 & {[}0.475, 0.567{]} & --- & 0.489
{[}0.300, 0.672{]} \\
Open − Closed（配对差） & +0.040 & {[}−0.022, +0.102{]} & --- & --- \\
\end{longtable}
}

调用量：探针共 1222 次模型调用（611 对 × 2 个条件），全部成功，0
次重试。另有 2 次正式运行前的冒烟调用，不计入上表。共约 56.6 万
token；按 \texttt{evals/\allowbreak{}run\_\allowbreak{}planner\_\allowbreak{}eval.\allowbreak{}py} 中的价格表估算，费用为
0.081 美元（off-peak）至 0.162 美元（peak）。

\subsection{解读}\label{ux89e3ux8bfb}

\begin{itemize}
\tightlist
\item
  \textbf{Open-book 的 AUC 显著高于 0.5}（CI 下界 0.526），而
  closed-book 和原始反转分数的 CI 都包含
  0.5。看到名称以后，模型确实从记忆中带进了与标签相关的信号。模型给出的高概率几乎都落在教科书级的适应症上，例如四环素类--痤疮（0.97--0.98）、氟哌啶醇--亨廷顿病（0.85）。这是\textbf{记忆/泄漏，不是发现}：模型没有使用任何表达数据。
\item
  \textbf{泄漏幅度比预期小}。open 与 closed 的配对差为 +0.04，95\% CI
  跨过 0（{[}−0.02, 0.10{]}）。原因是 TRANSCRIPT
  的很多"阳性"并不是常识性的适应症（例如环丙沙星--Wilson
  病、克拉霉素--Turner 综合征），模型对这些都给出约 0.02
  的低概率。反过来，一些被标为"未知"的配对其实是真实的临床适应症（辛伐他汀--急性心肌梗死、螺内酯--充血性心衰、氨苄西林--链球菌扁桃体炎，概率均为
  0.97）。模型的记忆和数据集标签只是部分重合。
\item
  \textbf{对 strict
  模式的意义}：即使在这个标签噪声很大的数据集上，名称可见仍然带来了可检测的、与表达信号无关的判别力（open-book
  的 CI 排除 0.5，closed-book 与原始反转分数的 CI 不排除 0.5）；但 open
  与 closed
  的直接配对差在本样本中还没有达到统计显著。在标签更接近教科书适应症的数据集上，这种泄漏只会更大。因此，凡是报告"LLM
  辅助药物重定位"的基准，都必须在 LLM 看不到身份信息的 strict
  模式下运行；open-book 的高分不能当作方法有效的证据。
\item
  Closed-book 的 AUC（0.527）和原始反转分数（0.519）相当，并且只输出 6
  种不同的概率值。模型基本上只是把反转数字单调地映射成很低的概率，没有额外的判别能力。
\end{itemize}

\subsection{局限}\label{ux5c40ux9650}

\begin{itemize}
\tightlist
\item
  \textbf{未知≠阴性}：0
  表示未记录，不表示无效。上面列举的真实适应症被标为 0，会直接压低
  open-book 的 AUC，所以 AUC
  是"与该数据集标签的一致性"，不是真实准确率。
\item
  \textbf{−1 只有 11 个}，相关的 AUC 只作描述。
\item
  \textbf{映射质量}：MedGen 的 title 有时是遗传易感性或亚型的条目（如
  ``Dermatitis, atopic, 2''、``Influenza, severe, susceptibility
  to''），这会削弱 open-book 的识别能力，所以泄漏量可能被低估。3
  个药物/疾病 ID 未映射。
\item
  只测了一个模型、一种问法，而且
  \texttt{temperature=0}；换成更强的模型或 few-shot 提问，泄漏可能更大。
\item
  阳性从 399 个中随机抽取了 300
  个；未知配对是均匀随机抽取，没有按药物或疾病匹配。
\item
  Closed-book 的提示里给了约 0.4\%
  的基准率，这会把输出概率整体压低，但不影响 AUC 的排序比较。
\end{itemize}

\subsection{复现}\label{ux590dux73b0}

\begin{Verbatim}[breaklines=true,breakanywhere=true]
python -m evals.contamination_probe map      # 映射（缓存后不再联网）
python -m evals.contamination_probe sample   # 冻结样本；哈希不同则拒绝覆盖
python -m evals.contamination_probe run      # 需要 .env 中的 DEEPSEEK_API_KEY；按 checkpoint 断点续跑
python -m evals.contamination_probe analyze
\end{Verbatim}

\end{document}
